% Sección aditiva. No modifica la edición de 211 páginas.
% Procedencia y comprobaciones: documentos internos de este mismo directorio.
\section{Raffinement nonadique, drapeaux positifs et compatibilité des moments}
\label{sec:dim-cofinal}

La longueur d'une représentation finie détermine le nombre de mineurs pouvant être formés avec ses colonnes. Le raffinement conserve les colonnes héritées et en ajoute d'autres dont la provenance est connue. La dimension d'une réalisation matricielle peut donc croître sans remplacer le registre qui fixe ses séparations. La construction suivante réunit cette croissance avec deux notions précises de compatibilité : la restriction des évaluations aux marques antérieures et la moyenne des observables sur les cellules descendantes. La première préserve les drapeaux polynomiaux ; la seconde conserve une mesure et sépare l'incrément d'information en détails orthogonaux.

La donnée de départ est la sortie dodécaphasique déjà produite par APP--TRIT--TPK. Les deux feuillets d'APP conservent leurs divisions complètes ; TRIT oriente le régime et l'opérateur \(U_t=\operatorname{Upd}_t\circ\operatorname{Tra}_t\circ
\operatorname{Sel}_t\) transporte l'état enrichi. Le prolongement \(w_6\to w_{12}\to w_{18}\to w_{24}\to w_{30}\to R_{36}\to G_9\) conserve le retour de phase et l'avancement de mémoire. La structure discrète conjointe du continu conserve corrélativement les représentations de cet état, avec ses feuillets, son orientation, sa mémoire et ses frontières. De ce même état proviennent les représentations régionales et le registre signé dont l'inversion donne
\[
 K=(234,543,140,729,659,824,621,58,914,794,146,601),
 \qquad Q=\sum_{j=1}^{12}K_j=6263.
\]
Cette sortie est utilisée avec ses marques et son orientation ; les matrices de
moments sont des réalisations ultérieures. L’indice de raffinement introduit
ci-dessous compte les subdivisions du lecteur géométrique. Son identification
à une durée physique ou à un nombre de microtransitions du TPK
exige de conserver l’horloge de cette autre réalisation.

\subsection{Partitions marquées et croissance cofinale}

Écrivons \(d_j=K_j/Q\), \(a_0=0\) et
\(a_j=\sum_{h=1}^j d_h\). Pour \(r\geq0\), nous définissons
\[
 L_r=12\,9^r,\qquad N_r=L_r+1,\qquad
 t_{r,(j-1)9^r+c}=a_{j-1}+\frac{c}{9^r}d_j,
 \quad 0\leq c\leq9^r.
\]
Les extrémités communes ne sont comptées qu'une fois. Les \(L_r\) cellules sont indexées par \((j,c)\), avec \(1\leq j\leq12\) et \(0\leq c<9^r\), et ont pour longueur \(\ell_{r,j,c}=d_j/9^r\). L'écriture de \(c\) en \(r\) chiffres nonadiques est également conservée, y compris les zéros initiaux. Un successeur a pour indice \(9c+h\), avec \(0\leq h\leq8\) ; la division euclidienne restitue son prédécesseur et le dernier chiffre. Les extrémités partagées retiennent leur marque de frontière ; la convention d'intervalles semi-ouverts ne fait qu'organiser les cellules et ne supprime aucune marque.

\begin{proposition}[Cofinalité et récupération du registre]
\label{prop:dim-grid}
Les marques satisfont
\[
 0=t_{r,0}<t_{r,1}<\cdots<t_{r,L_r}=1,
 \qquad t_{r+1,9i}=t_{r,i}.
\]
Le maximum des longueurs est
\(h_r=(\max_jK_j)/(Q9^r)\), qui tend vers zéro. Pour toute paire
d’entiers \(k\geq1\), \(m\geq1\), il existe \(r_0\) tel que
\(k+m\leq N_r\) pour \(r\geq r_0\). Le registre se retrouve à
tout niveau par
\[
 K_j=Q\sum_{c=0}^{9^r-1}\ell_{r,j,c}.
\]
\end{proposition}
\begin{proof}
Chaque différence consécutive est \(d_j/9^r>0\). En remplaçant \(c\) par \(9c\), la formule du niveau suivant donne exactement la marque antérieure. Les neuf longueurs des successeurs ont pour somme la longueur du prédécesseur ; l'itération prouve la récupération. Enfin, \(N_r=12\,9^r+1\) croît sans borne. Un choix explicite est tout \(r_0\geq0\) tel que \(9^{r_0}\geq(k+m-1)/12\). L'argument entier conserve \(Q\), l'indice de bloc et les chiffres de descendance.
\end{proof}

La normalisation détermine les rapports \(K_j/Q\). La récupération des entiers utilise en outre la charge \(Q\) conservée par le registre. On distingue ainsi la géométrie normalisée et l'information d'échelle accompagnant sa représentation.

\subsection{Drapeaux d'évaluations et réalisations en toute dimension finie}

Pour \(1\leq q\leq N_r\), soit
\[
 V_r^{(q)}=(t_{r,i}^{\,a})_{0\leq a<q,\,0\leq i<N_r},
 \qquad F_r^q=\operatorname{fila}V_r^{(q)}.
\]
La puissance d’exposant zéro vaut un également au nœud zéro.

\begin{theorem}[Drapeau positif et composition de rang arbitraire]
\label{thm:dim-all-ranks}
Les espaces \(F_r^q\) forment un drapeau \(F_r^1\subset F_r^2\subset\cdots\subset F_r^{N_r}\), avec \(\dim F_r^q=q\), et chaque matrice \(V_r^{(q)}\) a tous ses mineurs maximaux ordonnés strictement positifs. Pour \(k\geq1\), \(m\geq1\), \(k+m\leq N_r\), les matrices
\[
 C_r=V_r^{(k)},\quad Z_r=V_r^{(k+m)},\quad
 Y_r=C_rZ_r^{\mathsf T}
\]
satisfont
\[
 (Y_r)_{ab}=\sum_{i=0}^{L_r}t_{r,i}^{a+b},\qquad
 \operatorname{rank}Y_r=k,\qquad
 [Y_r]\in\mathcal A_{N_r,k,m}(Z_r).
\]
La construction atteint toute paire finie \((k,m)\) par un niveau
suffisamment profond de la même partition marquée.
\end{theorem}
\begin{proof}
Pour \(i_1<\cdots<i_q\), le mineur correspondant est
\[
 \det(t_{r,i_b}^{a-1})_{a,b=1}^{q}
   =\prod_{a<b}(t_{r,i_b}-t_{r,i_a})>0.
\]
Cela prouve le rang, la positivité et les inclusions du drapeau. Le bloc
initial \(k\times k\) de \(Y_r\) est le gramien des évaluations
de \(1,t,\ldots,t^{k-1}\). Pour \(x\neq0\),
\[
 x^{\mathsf T}(Y_r)_{[0,k-1],[0,k-1]}x
 =\sum_i\left(\sum_{a=0}^{k-1}x_at_{r,i}^{a}\right)^2>0,
\]
car un polynôme non nul de degré inférieur à \(k\) ne s’annule pas
aux \(N_r\geq k+m\) nœuds distincts. Le bloc est inversible,
et \(Y_r\), qui a \(k\) lignes, est de rang \(k\).
L’appartenance indiquée est l’application de la matrice externe positive
\(Z_r\) au point positif représenté par \(C_r\).
\end{proof}

On emploie la convention mathématique
\[
 \mathcal A_{N,k,m}(Z)=
 \{[CZ^{\mathsf T}]:[C]\in\operatorname{Gr}_{\geq0}(k,N)\}.
\]
La définition pour \(m\) général peut être consultée dans S.~N.~Karp, L.~K.~Williams et Y.~X.~Zhang,
\href{https://people.math.harvard.edu/~williams/papers/amplituhedron-plus.pdf}
{\emph{Decompositions of amplituhedra}}, définition~1.1.
Ici, la donnée externe provient déjà de la partition exprimée par HMT. La conclusion construit une collection déterminée de points et de drapeaux ; les dimensions de la grassmannienne ambiante et d'un atlas de son image sont des objets supplémentaires aux rangs démontrés.

Soit \(E_r:\mathbb R^{N_{r+1}}\to\mathbb R^{N_r}\) la restriction aux positions \(9i\). Pour tout \(q\leq N_r\),
\[
 E_r\bigl((V_{r+1}^{(q)})^{\mathsf T}x\bigr)
       =(V_r^{(q)})^{\mathsf T}x.
\]
Par conséquent, \(E_r:F_{r+1}^q\to F_r^q\) est un isomorphisme sur ces espaces d'évaluations : ses deux réalisations identifient le même polynôme de degré inférieur à \(q\). La composition des restrictions coïncide avec la restriction directe. Les mineurs sur les indices hérités conservent exactement leur valeur.

\begin{proposition}[Récupération à partir d’un drapeau calibré]
Si l’on conserve l’ordre des colonnes, les marques extrêmes et la
calibration \(t_{r,0}=0\), \(t_{r,L_r}=1\), le plan \(F_r^2\)
détermine les nœuds et les séparations. Avec \(Q\) et les
marques de descendance, il détermine \(K\).
\end{proposition}
\begin{proof}
Le plan contient le vecteur constant et le vecteur des nœuds. Toute
autre coordonnée qui, avec le vecteur constant, engendre le même
plan est \(a+bt\), avec \(b\neq0\). Les deux calibrages fixent
\(a=0\), \(b=1\). Les différences consécutives récupèrent les
longueurs, et la proposition~\ref{prop:dim-grid} récupère le registre.
\end{proof}

\subsection{Corps convexes de rang un et croissance du domaine}

Pour \(m\geq1\) et \(m+1\leq N_r\), nous définissons
\[
 P_r^{(m)}=\operatorname{conv}
 \{(t_{r,i},t_{r,i}^2,\ldots,t_{r,i}^m):0\leq i\leq L_r\}.
\]
La normalisation d'une ligne positive de \(C\) en coefficients barycentriques prouve \(\mathcal A_{N_r,1,m}(Z_r)=P_r^{(m)}\). Le déterminant de Vandermonde fournit \(m+1\) points affinement indépendants, de sorte que le corps a pour dimension \(m\).

\begin{proposition}[Recouvrement simplicial en dimension finie et limite convexe]
\label{prop:dim-polytopes}
Chaque \(P_r^{(m)}\) admet une triangulation par ses sommets, de
degré un sur l’intérieur de chaque cellule simpliciale. La somme de ses
formes canoniques a ses résidus intérieurs annulés. De plus,
\[
 P_r^{(m)}\subseteq P_{r+1}^{(m)},\qquad
 \overline{\bigcup_{r\geq r_0}P_r^{(m)}}=
 \operatorname{conv}\{(t,t^2,\ldots,t^m):0\leq t\leq1\}.
\]
\end{proposition}
\begin{proof}
On ordonne les sommets selon leurs marques. Dans chaque face, on choisit son plus petit
sommet ; on triangule récursivement les facettes qui ne le contiennent pas
et on forme les cônes correspondants. La récurrence sur la dimension
prouve la couverture et la disjonction des intérieurs, et la même règle sur les faces
communes garantit la compatibilité. Dans chaque simplexe orienté
\([v_0,\ldots,v_m]\), de matrice
\(A=(v_1-v_0,\ldots,v_m-v_0)\) à déterminant positif et de
coordonnées barycentriques \(\lambda_0,\ldots,\lambda_m\),
l’application barycentrique est affine et inversible. Sa forme est
\[
 \Omega_{[v_0,\ldots,v_m]}=
 \frac{dx_1\wedge\cdots\wedge dx_m}
 {\det A\prod_{j=0}^{m}\lambda_j}.
\]
La restriction résiduelle à chaque face est sa forme orientée. Deux simplexes adjacents induisent des orientations opposées sur leur face commune, de sorte que les résidus intérieurs s'annulent. C'est la même opération de triangulation et de résidu appliquée dans la réalisation de rang un précédente, le nombre de sommets étant maintenant augmenté par le raffinement.

L’inclusion des nœuds hérités prouve l’inclusion des corps. Le maillage
tend vers zéro, de sorte que les courbes échantillonnées sont denses dans la courbe
compacte des moments. Toute combinaison convexe finie de points de la
courbe s’approche par des combinaisons de nœuds d’un niveau commun.
Réciproquement, tous les corps sont contenus dans l’enveloppe
convexe compacte de la courbe. Les deux inclusions prouvent l’égalité.
\end{proof}

L'annulation des résidus opère sur une triangulation ou une subdivision d'un corps fixe. Le passage de \(P_r^{(m)}\) à \(P_{r+1}^{(m)}\) ajoute des sommets et peut agrandir le corps. Leurs formes appartiennent à ces domaines respectifs. La limite convexe précédente spécifie la convergence des corps, sans imposer d'identité entre leurs formes.

\subsection{Mesure du calendrier et convergence des moments nodaux}

Soit \(F_K:[0,1]\to[0,1]\) la fonction croissante affine par morceaux
qui envoie \((j-1)/12\) sur \(a_{j-1}\) et \(j/12\) sur \(a_j\).
Sur l’intervalle correspondant, sa pente est \(12d_j>0\).
La mesure du calendrier transporté est
\[
 \mu_K=(F_K)_*(ds),\qquad
 \frac{d\mu_K}{dt}=\frac1{12d_j}\quad
 (a_{j-1}<t<a_j).
\]
Chaque intervalle initial a pour masse \(1/12\), et chaque cellule de niveau
\(r\) a pour masse \(1/L_r\). Ce choix attribue le même poids
aux positions du calendrier avant de réaliser leurs longueurs.
La mesure de longueur \(dt\) est un autre lecteur ; les deux coïncident lorsque
toutes les séparations initiales sont égales.

Ses moments sont rationnels et déterminés par \(K\) :
\begin{equation}
 M_p=\int_0^1t^p\,d\mu_K(t)
   =\frac1{12(p+1)}\sum_{j=1}^{12}
       \frac{a_j^{p+1}-a_{j-1}^{p+1}}{d_j},\qquad p\geq0.
 \label{eq:dim-moments}
\end{equation}
\begin{proposition}[Récupération par moments et quantiles]
\label{prop:dim-moment-recovery}
La suite complète \((M_p)_{p\geq0}\) détermine \(\mu_K\). Si \(H_K(t)=\mu_K([0,t])\) est sa fonction de répartition, alors
\[
 a_j=H_K^{-1}(j/12),\qquad
 K_j=Q\bigl(H_K^{-1}(j/12)-H_K^{-1}((j-1)/12)\bigr).
\]
Par conséquent, les moments complets, la charge \(Q\) et l’ordre
du calendrier récupèrent le registre initial.
\end{proposition}
\begin{proof}
Deux probabilités sur \([0,1]\) ayant les mêmes moments ont les
mêmes intégrales pour tous les polynômes. L’approximation uniforme
des fonctions continues par des polynômes et la masse finie étendent cette
égalité à toutes les fonctions continues ; celles-ci déterminent la mesure
de Borel. La densité de \(\mu_K\) est positive et constante sur chaque
intervalle initial. Sa fonction de répartition est donc continue et
strictement croissante, et vérifie \(H_K(a_j)=j/12\). L’inversion
de cette égalité et la différence des extrémités prouvent les formules.
\end{proof}

La probabilité nodale correspondant au théorème \ref{thm:dim-all-ranks} est \(\nu_r=N_r^{-1}\sum_{i=0}^{L_r}\delta_{t_{r,i}}\).

\begin{theorem}[Limite de la représentation nodale]
\label{thm:dim-nodal-limit}
On a \(\nu_r\rightharpoonup\mu_K\). Si \(f\) est continue,
\[
 \left|\int f\,d\nu_r-\int f\,d\mu_K\right|
 \leq \omega_f(h_r)+\frac{\operatorname{osc}(f)}{N_r},
\]
où \(\omega_f\) est son module de continuité. Pour \(p\geq1\), l'erreur du moment d'ordre \(p\) est bornée par \(p h_r+1/N_r\). Pour chaque paire fixée \((k,m)\),
\[
 \frac1{N_r}Y_r\longrightarrow (M_{a+b})_{
     0\leq a<k,\,0\leq b<k+m}=:Y_\infty^{(k,m)},
 \qquad\operatorname{rank}Y_\infty^{(k,m)}=k.
\]
\end{theorem}
\begin{proof}
La mesure qui assigne une masse \(1/L_r\) à chaque extrémité gauche de
cellule approche son intégrale avec une erreur au plus égale à \(\omega_f(h_r)\).
En effet, chaque cellule a cette masse et un diamètre au plus égal à \(h_r\).
La mesure \(\nu_r\) est \(L_r/N_r\) fois cette mesure plus
\(\delta_1/N_r\), ce qui ajoute au plus
\(\operatorname{osc}(f)/N_r\). Pour \(f(t)=t^p\), la constante
de Lipschitz est \(p\). La convergence de chaque entrée découle de
ces bornes. Le bloc initial de la matrice limite est défini positif :
l’intégrale du carré d’un polynôme non nul est positive, car
\(\mu_K\) a une densité strictement positive sur tout intervalle
ouvert de \([0,1]\). La limite conserve ainsi le rang \(k\).
\end{proof}

Les moments nodaux varient exactement par l'incorporation de nœuds. Si \(\mathcal J_r\) est l'ensemble des nouvelles positions et \(z_q(t)=(1,t,\ldots,t^{q-1})^{\mathsf T}\), alors
\[
 V_{r+1}^{(q)}(V_{r+1}^{(q)})^{\mathsf T}
 =V_r^{(q)}(V_r^{(q)})^{\mathsf T}
  +\sum_{i\in\mathcal J_r}z_q(t_{r+1,i})z_q(t_{r+1,i})^{\mathsf T}.
\]
Chaque terme est positif semi-défini. Pour les matrices rectangulaires \(Y_r\), la même actualisation vaut avec \(z_kz_{k+m}^{\mathsf T}\). L'égalité exprime le changement de la mesure de comptage ; la normalisation et le théorème précédent décrivent sa limite. L'objet limite de rang \(k\) possède des coordonnées dans une grassmannienne finie, tandis que les données externes \(Z_r\) appartiennent aux différentes représentations de chaque niveau.

\subsection{Moyennes cellulaires, opérateurs compatibles et détail orthogonal}

Pour une cellule \(I_{r,i}=[b_{r,i},b_{r,i}+\ell_{r,i}]\), nous posons
\[
 a_{r,i}^{(p)}=
 \frac{(b_{r,i}+\ell_{r,i})^{p+1}-b_{r,i}^{p+1}}
 {(p+1)\ell_{r,i}},\qquad
 A_r^{(q)}=(a_{r,i}^{(p)})_{0\leq i<L_r,\,0\leq p<q}.
\]
Ce sont les moyennes des monômes par rapport à la probabilité conditionnelle
de \(\mu_K\) dans chaque cellule. Cette probabilité est uniforme dans la
cellule, bien que \(\mu_K\) ne soit pas la mesure de longueur globale.

Soit \(\mathcal H_r=\mathbb R^{L_r}\) muni du produit \(\langle x,y\rangle_r=L_r^{-1}\sum_i x_i y_i\). L'application \(J_r:\mathcal H_r\to\mathcal H_{r+1}\) répète la valeur d'un prédécesseur dans ses neuf successeurs. Son adjointe \(R_r=J_r^*\) prend la moyenne de ces neuf valeurs. On a \(R_rJ_r=I\).

\begin{theorem}[Compatibilité exacte et bilan des moments]
\label{thm:dim-cell-balance}
Pour tout \(p\geq0\),
\[
 R_r a_{r+1}^{(p)}=a_r^{(p)},\qquad
 \frac1{L_r}\sum_i a_{r,i}^{(p)}=M_p.
\]
Définissons
\[
 D_r^{(q)}=A_{r+1}^{(q)}-J_rA_r^{(q)},\quad
 G_r^{(q)}=\frac1{L_r}(A_r^{(q)})^{\mathsf T}A_r^{(q)}.
\]
Alors
\[
 R_rD_r^{(q)}=0,\qquad
 G_{r+1}^{(q)}=G_r^{(q)}+
       \frac1{L_{r+1}}(D_r^{(q)})^{\mathsf T}D_r^{(q)},
\]
et
\[
 G_r^{(q)}\longrightarrow H_K^{(q)}=(M_{a+b})_{0\leq a,b<q}.
\]
Les détails et le niveau initial reconstruisent tous les niveaux. Pour
\(q\leq L_r\), la matrice \(A_r^{(q)}\) a rang \(q\) et
tous les mineurs ordonnés d’ordre \(q\) sont positifs.
\end{theorem}
\begin{proof}
Les masses des neuf successeurs sont égales et leurs intervalles partitionnent l'intervalle prédécesseur. L'additivité de l'intégrale donne la première identité ; la somme sur toutes les cellules donne la seconde. La décomposition \(A_{r+1}=J_rA_r+D_r\) est orthogonale dans chaque colonne puisque \(R_rD_r=0\). En développant son Gram, les termes croisés s'annulent, et \(J_r\) est une isométrie. Cela prouve le bilan matriciel.

L'observable constant par cellules qui prend la valeur \(a_{r,i}^{(p)}\) approche uniformément \(t^p\), avec une erreur ne dépassant pas \(p h_r\) pour \(p\geq1\). Il converge donc dans \(L^2(\mu_K)\), et les produits scalaires convergent vers \(M_{a+b}\). La reconstruction s'obtient en itérant \(A_{r+1}=J_rA_r+D_r\).

Pour prouver les mineurs, choisissons \(q\) cellules ordonnées. Par
multilinéarité et Fubini, le déterminant de leurs moyennes est la moyenne
sur le produit de ces cellules de
\(\prod_{a<b}(x_b-x_a)\). Cet intégrande est positif à
l’intérieur du produit car les cellules sont disjointes et ordonnées.
Leurs frontières sont de mesure nulle. La moyenne est strictement positive.
\end{proof}

L'identité se prolonge à la paire de degrés \(k\) et \(k+m\) :
\[
 \overline Y_r=\frac1{L_r}(A_r^{(k)})^{\mathsf T}A_r^{(k+m)},
 \qquad
 \overline Y_{r+1}-\overline Y_r
 =\frac1{L_{r+1}}(D_r^{(k)})^{\mathsf T}D_r^{(k+m)}.
\]
Pour \(k+m\leq L_r\), les matrices de moyennes définissent une
seconde réalisation amplituédrique positive, à \(L_r\) colonnes,
et \(\overline Y_r\) est de rang \(k\). Ce sont des moyennes de cellules,
tandis que \(Y_r\) utilise des évaluations en \(N_r=L_r+1\) nœuds. Les deux
réalisations convergent vers \(Y_\infty^{(k,m)}\) ; leur coïncidence
à la limite découle des démonstrations précédentes.

La version opératorielle est également explicite. Pour une fonction continue \(f\), soit \(\Psi_r(f)\) l'opérateur diagonal de ses moyennes conditionnelles dans \(\mathcal H_r\). Il est linéaire, positif et unital, et
\[
 R_r\Psi_{r+1}(f)J_r=\Psi_r(f).
\]
Cette égalité se prouve en appliquant les deux membres à chaque vecteur de base.
Pour \(X(t)=t\), les deux premiers opérateurs déterminent, cellule par
cellule,
\[
 \Psi_r(X)_i=b_{r,i}+\frac{\ell_{r,i}}2,\qquad
 \Psi_r(X^2)_i-\Psi_r(X)_i^2=\frac{\ell_{r,i}^2}{12}.
\]
Si \(v_{r,i}\) désigne cette variance, la formule \(\ell_{r,i}=\sqrt{12v_{r,i}}\) récupère la longueur ; la première identité récupère l'extrémité initiale. Avec l'orientation, les marques et \(Q\), on récupère à nouveau \(K\). La moyenne est une compression positive : la variance écrite quantifie pourquoi \(\Psi_r(X^2)\) et \(\Psi_r(X)^2\) sont des opérations distinctes.

\subsection{Transport de la mesure et portée de la croissance dimensionnelle}

Deux listes positives de séparations \(d\) et \(d'\), avec les
mêmes marques, déterminent un homéomorphisme croissant affine par intervalles
\(F=F_{K'}\circ F_K^{-1}\). Il envoie chaque cellule sur sa cellule homologue
et vérifie \(F_*\mu_K=\mu_{K'}\). C’est pourquoi
\[
 \mathcal U_F:L^2(\mu_K)\longrightarrow L^2(\mu_{K'}),
 \qquad\mathcal U_F f=f\circ F^{-1}
\]
est une isométrie surjective. Si \(P_r\) et \(P'_r\) sont les projections sur les fonctions constantes par cellules, alors \(P'_r\mathcal U_F=\mathcal U_FP_r\) : les moyennes conditionnelles sont transportées parce que les masses des cellules correspondantes coïncident. Cela s'applique en particulier au flot positif de séparations dirigé par \(K\), en conservant la direction de chaque bloc dans ses successeurs. Les moments de la coordonnée spatiale peuvent changer ; l'isométrie transporte l'observable avec la mesure et ses projections.

La croissance démontrée est cofinale en toute dimension externe finie. Elle conserve un système de drapeaux positifs, une collection explicite de points amplituèdraux, des corps convexes complets en rang un et une mesure dont les opérateurs de moyenne sont exactement compatibles. La mémoire des détails permet d'inverser la compression entre niveaux. Les preuves de couverture et de résidus de rang un conservent leur domaine ; pour les rangs supérieurs, le présent résultat fixe les matrices, la positivité, la récupération, la compatibilité et la limite avec des types explicites. Cette séparation situe la portée de chaque construction dans le même développement sans confondre croissance de dimension, paramétrisation d'un point et classification de toutes les cellules d'une image.
